\documentclass[12pt,a4paper]{article}

% ══════════════════════════════════════════════════════════════
% PACKAGES
% ══════════════════════════════════════════════════════════════
\usepackage[utf8]{inputenc}
\usepackage[T1]{fontenc}
\usepackage{lmodern}
% headheight fix — must precede fancyhdr
\setlength{\headheight}{14pt}
\addtolength{\topmargin}{-2pt}
\usepackage[margin=2.5cm]{geometry}
\usepackage{amsmath,amssymb,amsthm}
\usepackage{mathtools}
\usepackage{graphicx}
\usepackage{xcolor}
\usepackage{mdframed}
\usepackage{booktabs}
\usepackage{array}
\usepackage{multirow}
\usepackage{longtable}
\usepackage{enumitem}
\usepackage{authblk}
\usepackage{fancyhdr}
\usepackage{titlesec}
\usepackage{parskip}
\usepackage{soul}
\usepackage{tcolorbox}
\tcbuselibrary{skins,breakable}
% hyperref LAST — texorpdfstring used in headers to avoid
% PDF string warnings
\usepackage{hyperref}

% ══════════════════════════════════════════════════════════════
% COLOURS
% ══════════════════════════════════════════════════════════════
\definecolor{headerblue}{RGB}{20,60,140}
\definecolor{claimbox}{RGB}{20,60,140}
\definecolor{warnbox}{RGB}{160,30,30}
\definecolor{lightgray}{RGB}{245,245,245}
\definecolor{midgray}{RGB}{200,200,200}
\definecolor{darkgray}{RGB}{80,80,80}
\definecolor{forestgreen}{RGB}{0,120,60}
\definecolor{amber}{RGB}{180,100,0}

% ════════════════════════════════════════���═════════════════════
% SECTION FORMAT
% ══════════════════════════════════════════════════════════════
\titleformat{\section}
  {\large\bfseries\color{headerblue}}
  {\thesection.}{0.6em}{}
  [\vspace{-0.3em}\textcolor{headerblue}
   {\rule{\linewidth}{0.8pt}}]

\titleformat{\subsection}
  {\normalsize\bfseries\color{headerblue}}
  {\thesubsection.}{0.5em}{}

\titleformat{\subsubsection}
  {\normalsize\bfseries\color{darkgray}}
  {\thesubsubsection.}{0.5em}{}

% ══════════════════════════════════════════════════════════════
% HEADER / FOOTER
% texorpdfstring prevents hyperref from trying to encode
% \cdot into the PDF string for the bookmark
% ══════════════════════════════════════════════════════════════
\pagestyle{fancy}
\fancyhf{}
\fancyhead[L]{\small\color{darkgray}%
  Lawson --- HD Attractor Geometry}
\fancyhead[R]{\small\color{darkgray}%
  OrganismCore \texorpdfstring{$\cdot$}{.} 2026-03-09}
\fancyfoot[C]{\small\thepage}
\renewcommand{\headrulewidth}{0.4pt}

% ══════════════════════════════════════════════════════════════
% CUSTOM ENVIRONMENTS
% ══════════════════════════════════════════════════════════════

\newenvironment{claimenv}{%
  \begin{tcolorbox}[
    colback=claimbox,
    colframe=claimbox,
    coltext=white,
    boxrule=0pt,
    arc=3pt,
    left=8pt, right=8pt,
    top=6pt, bottom=6pt,
    breakable
  ]
  \bfseries\centering
}{%
  \end{tcolorbox}
}

\newenvironment{warnenv}{%
  \begin{tcolorbox}[
    colback=red!5,
    colframe=warnbox,
    boxrule=1pt,
    arc=3pt,
    left=8pt, right=8pt,
    top=6pt, bottom=6pt,
    breakable
  ]
}{%
  \end{tcolorbox}
}

\newenvironment{derivenv}{%
  \begin{tcolorbox}[
    colback=lightgray,
    colframe=midgray,
    boxrule=0.5pt,
    arc=3pt,
    left=8pt, right=8pt,
    top=6pt, bottom=6pt,
    breakable
  ]
}{%
  \end{tcolorbox}
}

\newenvironment{predenv}[1]{%
  \begin{tcolorbox}[
    colback=blue!4,
    colframe=headerblue,
    boxrule=1pt,
    arc=3pt,
    title={\bfseries\color{white}#1},
    colbacktitle=headerblue,
    left=8pt, right=8pt,
    top=6pt, bottom=6pt,
    breakable
  ]
}{%
  \end{tcolorbox}
}

% ── Theorem environments ──────────────────────────────────────
\newtheorem{theorem}{Theorem}
\newtheorem{axiom}{Axiom}
\newtheorem{corollary}{Corollary}
\newtheorem{lemma}{Lemma}
\newtheorem{definition}{Definition}
\newtheorem{prediction}{Prediction}

% ══════════════════════════════════════════════════════════════
% HYPERREF SETUP
% ═════════════════════════════════════════════════════��════════
\hypersetup{
  colorlinks  = true,
  linkcolor   = headerblue,
  citecolor   = headerblue,
  urlcolor    = headerblue,
  pdftitle    = {Huntington's Disease as Attractor Geometry},
  pdfauthor   = {Eric Robert Lawson},
  pdfsubject  = {Identity Axis Framework --- HD Application},
  pdfkeywords = {Huntington's disease, Waddington landscape,
    attractor geometry, epigenetic, HDAC, CBP, medium
    spiny neurons, neurodegeneration, identity axis,
    OrganismCore}
}

% ══════════════════════════════════════════════════════════════
% CONVENIENCE MACROS
% ══════════════════════════════════════════════════════════════
\newcommand{\Rratio}{R}
\newcommand{\RHD}{R_{\mathrm{HD}}}
\newcommand{\IA}{\mathcal{A}}
\newcommand{\Hub}{\mathcal{H}}

% ══════════════════════════════════════════════════════════════
\begin{document}
% ══════════════════════════════════════════════════════════════

% ── TITLE BLOCK ───────────────────────────────────────────────
\begin{center}
  {\LARGE\bfseries\color{headerblue}
    Huntington's Disease as Attractor Geometry}\\[0.6em]

  {\large
    A Principles-First Derivation of the\\
    Stage-Stratified Intervention Protocol\\
    from the Identity Axis Framework:\\[0.4em]
    \textit{The Trigger, the Hub, the Identity Anchor,\\
    and the Geometric Case That an ``Incurable'' Disease\\
    Is Calculably Curable at Every Stage}}\\[1.0em]

  {\normalsize Eric Robert Lawson}\\[0.2em]
  {\small
    OrganismCore (Independent Research)\\
    \href{mailto:OrganismCore@proton.me}
      {OrganismCore@proton.me}\quad$\cdot$\quad
    ORCID:
    \href{https://orcid.org/0009-0002-0414-6544}
      {0009-0002-0414-6544}}\\[0.5em]

  {\small 2026-03-09}\\[0.7em]

  \rule{0.6\linewidth}{0.4pt}\\[0.5em]

  {\footnotesize
    \textit{Framework DOI:}\;
    \href{https://doi.org/10.5281/zenodo.18898788}
      {10.5281/zenodo.18898788}\\[0.2em]
    \textit{Repository:}\;
    \href{https://github.com/Eric-Robert-Lawson/attractor-oncology}
      {\texttt{github.com/Eric-Robert-Lawson/attractor-oncology}}
  }\\[0.5em]
  \rule{0.6\linewidth}{0.4pt}
\end{center}

\vspace{0.5em}

% ── CENTRAL CLAIM BOX ─────────────────────────────────────────
\begin{claimenv}
Huntington's Disease is not incurable.
It is a terminal false attractor in the Waddington
epigenetic landscape of medium spiny neurons, driven by a
known single-gene trigger, maintained by a known epigenetic
Hub, operating on a known Identity Anchor, with a
stage-stratified intervention protocol that is
\emph{geometrically complete} --- and every tool required
at every stage either already exists or is in active
clinical development as of 2026-03-09.
The disease is calculable. The intervention is derivable.
The cure is achievable, defined precisely and stratified
by stage.
\end{claimenv}

\vspace{1em}

% ══════════════════════════════════════════════════════════════
\begin{abstract}
\noindent
We apply the Identity Axis Framework --- a principles-first,
axiomatic model of disease as attractor failure in Waddington
epigenetic state space --- to Huntington's Disease (HD).
We derive the cell of origin (medium spiny neurons, MSNs),
the Identity Anchor (BCL11B\slash CBP\slash DARPP-32\slash
BDNF-TrkB), and the Convergence Hub (HDAC2/3 acting unopposed
at MSN identity loci due to CBP sequestration by mutant
huntingtin, mHTT). We distinguish the \emph{trigger} (mHTT)
from the \emph{Hub} (HDAC2/3 dominance) and show this
distinction requires a two-target intervention strategy that
the current clinical landscape has not yet assembled.

We identify three coexisting MSN populations at any disease
stage (protected, rescuable, terminal) and derive a
stage-stratified four-level intervention protocol: Level~0
(mHTT substrate removal), Level~1 (Hub inhibition via
HDAC2/3 inhibitors), Level~1A (Identity Anchor restoration
via BDNF-TrkB agonism and CBP rescue), and Level~2A (cell
replacement via iPSC-derived MSN transplantation).

We show that: (i)~for pre-manifest carriers, the protocol
predicts \emph{prevention of disease onset} using existing
tools; (ii)~for early-manifest patients, near-halt of
progression is achievable with the Level~0\,+\,Level~1
combination; (iii)~for late-manifest patients, a defined
engineering pathway to cure exists. The current clinical
ceiling --- 75\% slowing from AMT-130 alone (uniQure
Phase~1/2, September 2025) --- is the Level~0 ceiling; the
geometry predicts it rises to $\geq$90\% slowing with the
full combination protocol.

Six falsifiable predictions are stated for the record, all
timestamped 2026-03-09. Zero structural contradictions with
the Identity Axis Theorem are present.

\medskip\noindent\textbf{Keywords:}
Huntington's disease; Waddington landscape; attractor
geometry; medium spiny neuron; CBP; HDAC; Identity Axis;
epigenetic Hub; mHTT; BCL11B; DARPP-32; BDNF; cell therapy;
neurodegeneration; principles-first derivation; pre-manifest
prevention
\end{abstract}

\newpage
\tableofcontents
\newpage

% ══════════════════════════════════════════════════════════════
\section{Preamble: Why This Disease Is Structurally Special}
% ══════════════════════════════════════════════════════════════

Every other disease in the Identity Axis Framework required
the derivation to \emph{find} the Identity Anchor and Hub
from the geometry of the epigenetic landscape --- to
identify, without prior knowledge, which gene most
constitutes the cell's identity and which gene most maintains
the false attractor. Huntington's Disease hands you the
answer.

It is the only major disease in the table where:
\begin{enumerate}[label=(\alph*)]
  \item The causal gene is \textbf{fully known} --- a single
    gene, \textit{HTT}, with a measurable CAG repeat
    expansion that predicts age of onset with quantitative
    precision;
  \item \textbf{Pre-symptomatic detection} is possible
    decades before onset via genetic testing, with 100\%
    certainty;
  \item The molecular cascade is among the most precisely
    characterised in all of neurodegeneration;
  \item The cell of origin is \textbf{unambiguous} ---
    medium spiny neurons (MSNs) of the caudate and putamen;
  \item The disease follows a \textbf{predictable trajectory}
    from a known starting condition: CAG repeat length
    predicts age of onset.
\end{enumerate}

This is the framework's most constrained case. Most
constrained means most tractable. The geometry is not
obscured by heterogeneity, unknown origin cells, or
ambiguous Hubs. It is maximally clear.

\begin{warnenv}
\textbf{The central claim of this document:} HD is
geometrically curable --- in the same sense that drag force
is calculable: given the correct inputs (stage, geometry,
intervention level), the outcome is derivable before any
clinical data is examined. The tools for every level either
exist or are in late-stage clinical development. The
combination has not been assembled. Assembling it is not
scientifically impossible. It is logistically and
organisationally incomplete.
\end{warnenv}

% ════════════════════════════���═════════════════════════════════
\section{Framework Foundations}
% ══════════════════════════════════════════════════════════════

\subsection{The Identity Axis Theorem}

The Identity Axis Framework rests on five axioms~\cite{IAT2026}.

\begin{axiom}[Waddington Principle]
Every normal cell type occupies a stable attractor in an
epigenetic landscape. The attractor is defined by the gene
regulatory network that maintains the cell's identity.
\end{axiom}

\begin{axiom}[False Attractor]
Disease is a cell displaced from its normal attractor and
stabilised in a false attractor, maintained by a convergent
epigenetic mechanism.
\end{axiom}

\begin{axiom}[Identity Anchor]
Every cell of origin has a defining identity programme
$\IA$ --- the element whose activity most directly
constitutes what the cell is.
\end{axiom}

\begin{axiom}[Convergence Hub]
Every false attractor has a maintaining element $\Hub$
--- whose activity most directly stabilises the displaced
state.
\end{axiom}

\begin{axiom}[Identity Axis]
The ratio
\begin{equation}
  \Rratio \;=\; \frac{\IA}{\Hub}
  \label{eq:ratio}
\end{equation}
is a continuous, measurable coordinate in the Waddington
landscape. $\Rratio$ predicts attractor depth. Attractor
depth predicts clinical outcome. $\Hub$ is the primary drug
target.
\end{axiom}

\subsection{HD as a Special Case: Terminal False Attractor}

\begin{definition}[Terminal False Attractor]
A false attractor is \emph{terminal} if the cell displaced
into it cannot survive indefinitely and progresses
irreversibly toward apoptosis.
\end{definition}

\begin{definition}[Stable False Attractor]
A false attractor is \emph{stable} if the displaced cell
can survive and self-maintain indefinitely in the false
state.
\end{definition}

In cancer, the false attractor is stable --- the cancer cell
lives, the Hub maintains it, and the cell is recoverable
by Hub dissolution. In HD, the false attractor is
\textbf{terminal}: the MSN cannot sustain the displaced
state, attempts cell-cycle re-entry, and undergoes
apoptosis. Every dead MSN is a permanent loss.

\begin{corollary}[Intervention Window]
The fraction of rescuable MSNs decreases monotonically with
disease progression. The pre-manifest stage maximises the
rescuable fraction. The intervention window is finite.
\end{corollary}

% ══════════════════════════════════════════════════════════════
\section{The Geometric Mapping}
% ══════════════════════════════════════════════════════════════

\subsection{Cell of Origin}

MSNs comprise $\approx$95\% of striatal neurons.
Within this population, \emph{indirect-pathway} D2-type
MSNs (expressing DRD2 and PENK/enkephalin, projecting
to external globus pallidus) degenerate first.

This ordering is geometrically predicted: D2 MSNs have
lower baseline BDNF-TrkB signalling density than D1
MSNs~\cite{Oxford2024}, placing them closer to the identity
maintenance threshold at baseline. Cells with the thinnest
margin fall first. This is attractor depth ordering: the
shallowest attractor yields first under Hub pressure.

\subsection{The Identity Anchor \texorpdfstring{$(\IA)$}{(A)}}

The MSN Identity Anchor operates at three nested levels.

\subsubsection{Deep Identity Anchor: BCL11B\,/\,CTIP2}

BCL11B is the master transcription factor for MSN identity
specification. It controls differentiation of MSN progenitors
from the lateral ganglionic eminence, governs all downstream
MSN-selective gene expression, and represses alternative
neuronal fates~\cite{Arlotta2008}. BCL11B loss in
experimental models recapitulates core HD features:
DARPP-32 loss, MSN identity marker loss, and striatal
connectivity disruption~\cite{Desplats2008,Tang2012}.

\subsubsection{Operational Identity Anchor:
CBP-Mediated Chromatin Accessibility}

CREB-binding protein (CBP) is a histone acetyltransferase
(HAT) that maintains chromatin accessibility at BCL11B,
DARPP-32, BDNF, and PGC-1$\alpha$ loci via H3K27
and H3K9 acetylation. Without CBP HAT activity at these
loci the MSN Identity Anchor programme cannot be maintained
regardless of which transcription factors are present.

\subsubsection{Trophic Identity Anchor: BDNF\,/\,TrkB}

BDNF is produced in the cortex and transported anterogradely
to the striatum, where it activates TrkB receptors on MSNs.
Loss of TrkB signalling in indirect-pathway SPNs triggers
GSTO2 upregulation, leading to aberrant glutathione
metabolism, elevated striatal dopamine, and dopaminergic
neurotoxicity to MSNs~\cite{Oxford2024}. This cascade
precedes clinical onset.

\medskip
The Identity Anchor in ratio terms is operationally
measured as:
\begin{equation}
  \IA \;=\;
  \underbrace{
    \mathrm{H3K27ac} + \mathrm{H3K9ac}
  }_{\substack{
    \text{at BCL11B, DARPP-32,}\\
    \text{BDNF, PGC-1}\alpha\text{ loci}
  }}
  \label{eq:anchor}
\end{equation}

Clinical surrogate: DARPP-32 expression and CSF BDNF levels.

\subsection{The Convergence Hub
\texorpdfstring{$(\Hub)$}{(H)}}

\begin{warnenv}
\textbf{Critical distinction --- mHTT is not the Hub.}\\
mHTT is the \emph{trigger}: the upstream force that
initiates the cascade. The Hub is the element that
\emph{maintains} the false attractor once displacement has
begun. These are different molecular entities. Targeting
only the trigger is necessary but not sufficient.
\end{warnenv}

\noindent The Hub is: \textbf{HDAC2/3 acting unopposed at
MSN identity loci due to CBP sequestration by mHTT.}

The molecular cascade:

\begin{derivenv}
\[
\underbrace{
  \text{mHTT (polyQ)}
}_{\text{trigger}}
\;\xrightarrow{\text{binds CBP}}\;
\underbrace{
  \text{CBP sequestered}
}_{\text{in aggregates}}
\;\xrightarrow{\text{HAT absent}}\;
\underbrace{
  \text{HDAC2/3 unopposed}
}_{\text{at MSN loci}}
\]
\[
\xrightarrow{\hphantom{\text{HDAC2/3 unopposed}}}
\underbrace{
  \mathrm{H3K27ac}\downarrow\;/\;\mathrm{H3K9ac}\downarrow
}_{\text{BCL11B, DARPP-32, BDNF, PGC-1}\alpha}
\;\xrightarrow{\;\Rratio\to 0\;}\;
\underbrace{
  \text{cell-cycle re-entry}
}_{\text{terminal false attractor}}
\;\longrightarrow\;
\text{apoptosis}
\]
\end{derivenv}

\noindent\textbf{Literature confirmations:}
\begin{itemize}
  \item HDAC2/3 selective knockout in MSNs partially
    corrects HD transcriptional pathology~\cite{eLife2020};
  \item Vorinostat (SAHA), brain-penetrant pan-HDAC
    inhibitor, extends survival and improves motor
    behaviour in HD mouse models~\cite{HDACiReview2020};
  \item HDAC6-selective inhibitor CKD-504 reduces mHTT
    aggregation \emph{and} restores synaptic function in
    YAC128 HD mice~\cite{CKD504_2023};
  \item Selective HDAC3 inhibition rescues synaptic and
    transcriptional phenotypes in HD models~\cite{HDAC3HD}.
\end{itemize}

% ══════════════════════════════════════════════════════════════
\section{The Three MSN Populations}
% ══════════════════════════════════════════════════════════════

At any disease stage, three coexisting MSN populations are
derivable from the geometry. Their proportions shift
monotonically with progression (Table~\ref{tab:populations}).

\subsection{Population~1 --- Protected
(\texorpdfstring{$\Rratio$}{R} above threshold)}

$\Rratio$ is above the identity maintenance threshold.
DARPP-32 is reduced but not lost. The cell is stressed but
alive and identified.

\textbf{These cells can be fully protected.} Remove the
trigger; the cascade does not advance; $\Rratio$ stabilises;
the cell survives with identity intact.

\subsection{Population~2 --- Rescuable
(\texorpdfstring{$\Rratio$}{R} below threshold, cell alive)}

$\Rratio$ has fallen below threshold. HDAC2/3 dominance is
active. DARPP-32 and BCL11B are significantly reduced.
The cell approaches the terminal threshold but has not
crossed it. Chromatin at MSN identity loci is closed but
not irreversibly so.

\textbf{These cells can be rescued.} If Hub dominance is
removed AND the Identity Anchor signal is restored,
$\Rratio$ rises back above threshold. This is the
mechanistic basis for the 75\% slowing with AMT-130
alone~\cite{AMT130_2025}: Population~2 cells partially
rescued by trigger removal even without explicit Hub
inhibition. Level~1 Hub inhibition should rescue
significantly more.

\subsection{Population~3 --- Terminal}

$\Rratio$ has collapsed. Cell-cycle re-entry has occurred.
Apoptosis is underway or complete.

\textbf{These cells cannot be rescued by Levels~0--1A.}
The intervention is replacement.

\begin{table}[ht!]
\centering
\caption{MSN population distribution by disease stage
  (geometric estimate)}
\label{tab:populations}
\begin{tabular}{lcccc}
\toprule
\textbf{Stage} &
\textbf{Pop.~1 (\%)} &
\textbf{Pop.~2 (\%)} &
\textbf{Pop.~3 (\%)} \\
\midrule
Pre-manifest   & 95--99 & 1--5  & $\approx$0  \\
Early manifest & 60--75 & 20--35 & 5--10 \\
Mid-manifest   & 20--40 & 20--30 & 30--50 \\
Late manifest  & 5--15  & 5--10  & 75--90 \\
\bottomrule
\end{tabular}
\end{table}

% ══════════════════════════════════════════════════════════════
\section{The Four Intervention Levels}
% ══════════════════════════════════════════════════════════════

\subsection{Level~0 --- Substrate Removal}

\textbf{Target:} mHTT protein.\\
\textbf{Mechanism:} Reduce mHTT below the threshold at which
CBP sequestration is sufficient to collapse the Identity
Anchor programme.

\medskip\noindent\textbf{Available tools (as of 2026-03-09):}

\begin{itemize}
  \item \textbf{AMT-130} (uniQure) --- AAV5-delivered
    artificial miRNA; non-allele-selective; one-time
    intraparenchymal delivery to caudate and putamen;
    Phase~1/2 topline result (September~2025): 75\%
    slowing of disease progression at 36~months on cUHDRS
    vs.\ propensity-matched controls; NfL below baseline
    at 36~months; FDA Breakthrough Therapy and RMAT
    designations~\cite{AMT130_2025}. FDA requires a
    randomised Phase~3 trial (March~2026).

  \item \textbf{WVE-003} (Wave Life Sciences) ---
    allele-selective ASO targeting HTT SNP3 in
    heterozygous patients; Phase~1b/2a SELECT-HD result
    (2024): 46\% reduction in CSF mHTT at week~24,
    wild-type HTT preserved and increased; caudate atrophy
    slowing confirmed~\cite{WVE003_2024}.
    \emph{Geometrically preferred}: wild-type HTT serves
    essential cellular functions; preserving it while
    removing mHTT is the precise intervention.

  \item \textbf{Tominersen} (Roche/Ionis) ---
    non-allele-selective intrathecal ASO; GENERATION~HD2
    Phase~2 ongoing~\cite{GenHD2}.

  \item \textbf{SKY-0515} (Skyhawk) --- oral small
    molecule splice modulator; FALCON-HD Phase~2/3
    initiated January~2026~\cite{SKY0515}.
\end{itemize}

\textbf{Geometric assessment:}
Level~0 is necessary but not sufficient alone.
75\% slowing is the Level~0 ceiling. Hub inhibition is
required to raise it.

\subsection{Level~1 --- Hub Inhibition}

\textbf{Target:} HDAC2/3 acting unopposed at BCL11B,
DARPP-32, BDNF, and PGC-1$\alpha$ chromatin loci.\\
\textbf{Mechanism:} Restore H3K27ac and H3K9ac at MSN
Identity Anchor loci $\to$ chromatin reopens $\to$
$\Rratio$ rises above threshold $\to$ Population~2 rescues.

\medskip\noindent\textbf{Available tools:}

\begin{itemize}
  \item \textbf{Vorinostat (SAHA)} --- FDA-approved
    pan-HDAC inhibitor; brain penetrant; extends survival
    in HD mouse models; restores histone acetylation at
    HD-relevant loci~\cite{HDACiReview2020}.
    \emph{Not currently in HD clinical trials as a primary
    intervention. The geometry indicates it should be.}

  \item \textbf{HDAC2/3-selective inhibitors} ---
    preferred over pan-HDAC; HDAC2/3 knockout in MSNs
    shows partial correction of HD transcriptional
    pathology~\cite{eLife2020}; multiple preclinical
    programmes active.

  \item \textbf{CKD-504 (HDAC6-selective)} --- reduces
    mHTT aggregation AND restores synaptic function in
    YAC128 mice~\cite{CKD504_2023}. Addresses both the
    epigenetic Hub and the protein clearance deficit.
\end{itemize}

\textbf{The missing clinical combination:}
Level~0\,+\,Level~1 simultaneously has not been run in
any clinical trial. This is the highest-priority next HD
trial.

\subsection{Level~1A --- Identity Anchor Restoration}

\textbf{Target:} Directly restore CBP HAT activity and
BDNF-TrkB signalling at MSN identity loci.

\medskip\noindent\textbf{Tools:}

\begin{itemize}
  \item \textbf{CREB pathway agonism:} PDE4 inhibitors
    raise cAMP $\to$ PKA activation $\to$ CREB Ser133
    phosphorylation $\to$ CBP recruitment to CRE sites.
    Roflumilast (brain-penetrant PDE4i, FDA-approved for
    COPD) has not been tested in HD as a primary
    intervention.

  \item \textbf{TrkB agonism:} Restores the trophic
    Identity Anchor signal. The Oxford/Nature~2024
    paper~\cite{Oxford2024} establishes that TrkB loss
    in iSPNs triggers the GSTO2-hyperdopaminergia cascade
    mediating earliest HD pathology. TrkB restoration
    addresses this cascade directly. BDNF pathway
    biomarker study NCT04012411 is
    active~\cite{BDNF_NCT}.

  \item \textbf{dCas9-p300:} Targeted delivery of the
    CBP catalytic HAT domain to specific MSN identity
    loci. Exists in preclinical tools; not yet applied
    to HD.
\end{itemize}

\subsection{Level~2A --- Identity Rebuilding and
Cell Replacement}

\textbf{For Population~3 --- cells beyond rescue.}

\subsubsection{Approach~A: BCL11B Re-specification}

AAV-BCL11B delivery forces re-specification of remaining
MSNs toward the correct identity programme. Pre-condition:
Ki67-negative MSN status at delivery (not yet in
cell-cycle re-entry). Not yet tested in HD.

\subsubsection{Approach~B: iPSC-Derived MSN
Transplantation}

\begin{enumerate}[label=\arabic*.]
  \item Patient cells $\to$ iPSC reprogramming;
  \item CAG repeat correction by allele-specific
    CRISPR-Cas9 \emph{before} differentiation (required:
    uncorrected transplants develop the same cascade ---
    see Prediction~6);
  \item BCL11B-directed MSN differentiation: ${>}90\%$
    purity in 21~days confirmed~\cite{iPSC_MSN_2024};
  \item Stereotactic delivery to caudate and putamen.
\end{enumerate}

First-in-human iPSC-MSN transplantation trials for HD:
anticipated late 2020s. Active preclinical
programme~\cite{iPSC_MSN_2024}.

\subsubsection{Approach~C: In Vivo Chemical Reprogramming
(Horizon)}

Direct in vivo conversion of surviving striatal astrocytes
to MSN identity via small molecule cocktails and BCL11B
viral delivery. Proof of concept for
astrocyte-to-neuron reprogramming exists; HD-specific
application is a 2030s research horizon.

% ═══════════════════════════════════════════════════��══════════
\section{The Stage-Stratified Intervention Protocol}
% ══════════════════════════════════════════════════════════════

\subsection{Stage~0: Pre-Manifest}

\textbf{Population profile:}
Pop.~1 $\approx$95--99\%;
Pop.~2 $\approx$1--5\%;
Pop.~3 $\approx$0\%.

\textbf{Required levels:} 0\,+\,1\,+\,1A.

\textbf{Protocol:} allele-selective ASO (WVE-003 or
successor), intrathecal, chronic maintenance; PLUS
HDAC2/3-selective inhibitor, oral, daily, brain-penetrant;
PLUS TrkB agonist to maintain corticostriatal BDNF-TrkB
signalling and prevent GSTO2-mediated cascade onset.

\begin{predenv}{Geometric Prediction --- Stage 0: Prevention}
If mHTT is reduced below the CBP sequestration threshold,
AND HDAC2/3 is inhibited to prevent Hub dominance from
establishing, AND BDNF-TrkB is maintained to prevent the
early metabolic cascade, then $\Rratio = \IA/\Hub$ never
falls below threshold. The false attractor never forms.
Population~2 never materialises. Population~3 never
materialises. \textbf{The disease never manifests.}

This is the most important prediction in this document.
For a disease in which carrier status is known at birth,
pre-manifest prevention using existing tools in combination
means HD becomes the first inherited neurodegenerative
condition preventable from onset in every identified carrier.
\end{predenv}

\subsection{Stage~1: Early Manifest}

\textbf{Population profile:}
Pop.~1 $\approx$60--75\%;
Pop.~2 $\approx$20--35\%;
Pop.~3 $\approx$5--10\%.

\textbf{Required levels:} 0\,+\,1\,+\,1A.

\begin{predenv}{Geometric Prediction --- Stage 1: Near-Halt}
AMT-130 alone (Level~0): 75\% slowing confirmed.
The continuing 25\% is Population~2 engaged before mHTT
reduction, plus Population~3 (5--10\%).

Level~0\,+\,Level~1\,+\,Level~1A: Population~1 fully
protected; Population~2 substantially rescued;
Population~3 (5--10\%) unrescuable by this protocol.
\textbf{Predicted outcome: $\geq$90--95\% slowing.}
At early-manifest stage, 90--95\% slowing is functionally
equivalent to halt for most patients over a normal lifespan.
\end{predenv}

\subsection{Stage~2: Mid-Manifest}

\textbf{Population profile:}
Pop.~1 $\approx$20--40\%;
Pop.~2 $\approx$20--30\%;
Pop.~3 $\approx$30--50\%.

\textbf{Required levels:}
0\,+\,1\,+\,1A; consideration of Level~2A Approach~A
(BCL11B AAV) for boundary cells.

\begin{predenv}{Geometric Prediction --- Stage 2:
Stabilisation with Partial Recovery}
Populations~1 and~2 (40--70\% of remaining MSNs)
protected or rescued. Population~3 (30--50\%) lost.
Some functional improvement as rescued MSNs restore
correct circuit function. Not halt. Not reversal of
existing loss. Meaningful clinical stabilisation and
partial recovery.
\end{predenv}

\subsection{Stage~3: Late Manifest}

\textbf{Population profile:}
Pop.~1 $\approx$5--15\%;
Pop.~2 $\approx$5--10\%;
Pop.~3 $\approx$75--90\%.

\textbf{Required levels:}
0\,+\,1 for remaining populations; Level~2A Approach~B
(iPSC-MSN transplantation) for functional restoration.

\begin{predenv}{Geometric Prediction --- Stage 3:
Stabilisation + Cell-Replacement-Dependent Recovery}
Levels~0\,+\,1 protect the $\approx$15--25\% of MSNs
still accessible. Severely depleted striatal circuitry
cannot be restored by cell rescue alone. Level~2A is
required for functional restoration. Timeline: late 2020s
at earliest for first-in-human trials. Not ``never.''
``Not yet.'' The engineering pathway is fully defined.
\end{predenv}

% ══════════════════════════════════════════════════════════════
\section{The Biomarker}
% ══════════════════════════════════════════════════════════════

The framework derives a specific biomarker directly
analogous to the IHC protocol for cancer:

\begin{equation}
  \RHD \;=\;
  \frac{
    \text{CBP HAT activity at MSN identity loci}
  }{
    \text{HDAC2/3 occupancy at MSN identity loci}
  }
  \label{eq:RHD}
\end{equation}

Operationally:
\begin{align}
  \IA &= \mathrm{H3K27ac} + \mathrm{H3K9ac}
      \;\text{at BCL11B, DARPP-32, BDNF, PGC-1}\alpha
      \;\text{loci}
      \label{eq:A_op}\\
  \Hub &= \text{HDAC2/3 ChIP signal at the same loci}
      \label{eq:H_op}
\end{align}

Clinical surrogates for trial use are summarised in
Table~\ref{tab:biomarkers}.

\begin{table}[ht!]
\centering
\caption{Biomarker hierarchy for the proposed combination
  trial}
\label{tab:biomarkers}
\begin{tabular}{p{3.2cm}p{3.5cm}p{4.8cm}}
\toprule
\textbf{Biomarker} &
\textbf{Measures} &
\textbf{Status} \\
\midrule
CSF mHTT &
Level~0 adequacy &
Validated~\cite{AMT130_2025,WVE003_2024} \\[3pt]
NfL (CSF/blood) &
Neuronal stabilisation &
Validated~\cite{AMT130_2025} \\[3pt]
Caudate MRI vol. &
Striatal loss rate &
Validated~\cite{WVE003_2024} \\[3pt]
DARPP-32 (CSF) &
MSN identity integrity &
Proposed --- novel \\[3pt]
GSTO2 (CSF) &
Earliest Identity Anchor
collapse (pre-NfL) &
Proposed --- novel \\[3pt]
H3K27ac at DARPP-32 &
$\RHD$ directly &
Proposed; striatal
biopsy required \\
\bottomrule
\end{tabular}
\end{table}

\textbf{Novel prediction:} GSTO2 in CSF will be elevated
in pre-manifest HD carriers prior to clinical onset and
will decline in response to TrkB agonist therapy
\emph{before} DARPP-32 levels change. GSTO2 is a leading
biomarker of Identity Anchor collapse, preceding both NfL
elevation and detectable mHTT change-from-baseline.

% ══════════════════════════════════════════════════════════════
\section{The Complete Answer: Is HD Curable?}
% ══════════════════════════════════════════════════════════════

\subsection{Definition of Cure}

In the framework's precise language, \emph{cure} means:
$\Rratio = \IA/\Hub$ is restored and maintained above the
identity threshold in the target cell population for the
duration of the patient's life.
This definition has four stage-specific instantiations
(Table~\ref{tab:cure}).

\begin{table}[ht!]
\centering
\caption{Stage-stratified definition of cure and achievability}
\label{tab:cure}
\begin{tabular}{p{2.0cm}p{3.2cm}p{3.4cm}p{3.0cm}}
\toprule
\textbf{Stage} &
\textbf{Cure definition} &
\textbf{Achievability} &
\textbf{Missing piece} \\
\midrule
Pre-manifest &
Prevention: disease never manifests &
\textcolor{forestgreen}{\textbf{YES}} with existing tools
in combination &
Combination trial not yet run \\[4pt]
Early manifest &
Near-halt: $\geq$90\% slowing &
\textcolor{forestgreen}{\textbf{YES}} with combination
trial &
Combination trial not yet run \\[4pt]
Mid-manifest &
Stabilisation + partial recovery &
\textcolor{amber}{\textbf{SUBSTANTIALLY}}
with combination + BCL11B consideration &
BCL11B AAV preclinical only \\[4pt]
Late manifest &
Stabilisation + cell-replacement recovery &
\textcolor{amber}{\textbf{DEFINED PATHWAY}}
iPSC-MSN required &
iPSC-MSN first-in-human late 2020s \\
\bottomrule
\end{tabular}
\end{table}

\subsection{What the Field Has. What Is Missing.}

The current clinical landscape has assembled Level~0 with
precision. What the field has \emph{not} assembled:

\begin{enumerate}
  \item \textbf{Level~0\,+\,Level~1 combination trial.}
    Vorinostat is FDA-approved and brain-penetrant.
    Preclinical confirmation exists. The clinical
    combination has not been run. This is the next trial.

  \item \textbf{Level~1A: TrkB agonism as co-intervention.}
    The Oxford/Nature~2024 paper established the GSTO2
    mechanism as the earliest HD pathology. The
    therapeutic consequence has not entered clinical
    trials.

  \item \textbf{The pre-manifest prevention trial.}
    Pre-manifest CAG-confirmed carriers receiving
    Level~0\,+\,1\,+\,1A prophylactically. This could make
    HD the first inherited neurodegenerative disease
    preventable from onset in every identified carrier.

  \item \textbf{DARPP-32 and GSTO2 as primary trial
    biomarkers.} Neither has been deployed as a primary
    endpoint in any HD trial.
\end{enumerate}

\subsection{Is This Ambitious?}

\begin{itemize}
  \item Stage~0: \textbf{No.} Tools exist. Biology
    confirmed. Missing piece is a clinical decision.

  \item Stage~1: \textbf{Moderately.} All tools exist or
    are in development. The combination trial is
    logistically but not scientifically ambitious.

  \item Stages~2--3: \textbf{Genuinely.} BCL11B AAV and
    iPSC-MSN transplantation require tools not yet in
    clinical use. Science is sound. Engineering is
    defined. Timeline is real.
\end{itemize}

% ══════════════════════════════════════════════════════════════
\section{Six Falsifiable Predictions for the Record}
% ══════════════════════════════════════════════════════════════

All predictions are derived from the framework and
timestamped 2026-03-09.

\begin{prediction}[Combination superiority]
\label{pred:combo}
AMT-130 (or allele-selective ASO) + HDAC2/3-selective
inhibitor in early-manifest HD patients will demonstrate
statistically significant superiority over Level~0 alone
on cUHDRS and TFC at 24~months. Additional benefit will be
mediated by Population~2 MSN rescue evidenced by DARPP-32
recovery in striatal tissue.
\end{prediction}

\begin{prediction}[Pre-manifest prevention]
\label{pred:premanifest}
Pre-manifest HD carriers receiving Level~0\,+\,1\,+\,1A
prophylactically will show delayed or absent conversion to
manifest HD compared to Level~0 alone. Prevention rate
will correlate with residual mHTT level and H3K27ac
recovery at DARPP-32 loci.
\end{prediction}

\begin{prediction}[GSTO2 as leading biomarker]
\label{pred:gsto2}
GSTO2 in CSF will be elevated in pre-manifest HD carriers
prior to clinical onset and will decline in response to
TrkB agonist therapy before DARPP-32 levels change. GSTO2
is a leading biomarker of Identity Anchor collapse,
preceding both NfL elevation and detectable mHTT
change-from-baseline.
\end{prediction}

\begin{prediction}[HDAC6\,+\,HDAC2/3 additivity]
\label{pred:hdac_add}
HDAC6-selective inhibition will show additive benefit with
HDAC2/3-selective inhibition in HD because they address
different downstream consequences of Hub dominance:
HDAC2/3 at MSN identity chromatin (gene expression rescue)
and HDAC6 at tubulin deacetylation (mHTT aggregate
clearance via autophagy).
\end{prediction}

\begin{prediction}[BCL11B AAV efficacy conditionality]
\label{pred:bcl11b}
BCL11B AAV delivery to mid-manifest HD striatum will show
dose-dependent DARPP-32 recovery in transduced MSNs and
correlate with motor function improvement at 12~months.
Rescue will be limited to Ki67-negative MSNs at delivery.
The boundary condition is chromatin accessibility, not
BCL11B expression level.
\end{prediction}

\begin{prediction}[Allele-corrected iPSC-MSN superiority]
\label{pred:ipsc}
Allele-corrected (CRISPR CAG-edited) iPSC-MSNs will show
superior long-term engraftment compared to uncorrected
transplants. Uncorrected transplants develop the same
Hub-driven cascade and progressively lose MSN identity.
Allele correction is required for durable therapeutic
benefit.
\end{prediction}

% ══════════════════════════════════════════════════════════════
\section{Falsifiability}
% ══════════════════════════════════════════════════════════════

Following the governing framework
principle~\cite{FalsThm2026}:

\subsection{Type~B Failure (framework-invalidating)}

If the combination of (i)~mHTT reduction below CBP
sequestration threshold, (ii)~HDAC2/3 inhibition to
restore H3K27ac at MSN identity loci, and (iii)~BDNF-TrkB
trophic support applied in early-manifest or pre-manifest
HD patients does \emph{not} produce superior outcomes to
Level~0 alone, AND no structural explanation within the
epigenetic landscape geometry accounts for the failure,
\textbf{then the framework's application to HD is
invalidated.}

\subsection{Type~A Failures (informative, not invalidating)}

\begin{itemize}
  \item \textit{HDAC2/3 inhibition alone insufficient:}
    Identity Anchor has two components; both required.
    Map extends. Axioms intact.

  \item \textit{Combination works for D1 but not D2:}
    D2 MSNs have lower baseline BDNF-TrkB density; their
    threshold position differs. Map extends. Axioms intact.

  \item \textit{BCL11B AAV does not rescue Population~3
    boundary cells:}
    Terminal threshold was already crossed. BCL11B requires
    accessible chromatin. Pre-condition is Ki67-negativity.
    Map extends. Axioms intact.
\end{itemize}

% ══════════════════════════════════════════════════════════════
\section{Discussion}
% ══════════════════════════════════════════════════════════════

\subsection{Why the ``Incurable'' Designation Persists}

HD has been designated incurable not because the biological
mechanisms were unknown but because the field has been
organised around a single therapeutic strategy (HTT
lowering) without the geometric framework that explains
why single-strategy approaches have a ceiling, what that
ceiling is, and what is required to exceed it.

AMT-130's 75\% slowing is a landmark achievement. It is
also, from the geometry, the \emph{expected} ceiling for
Level~0 alone. The 25\% that continues to progress is the
Population~2 and~3 fraction that Level~0 cannot reach.
The geometry explains the ceiling and explains how to
raise it.

\subsection{The Pre-Manifest Window Is the
Highest-Value Target}

The most important clinical trial that has never been run
in HD is the pre-manifest combination trial. The disease
is genetically predictable. Carriers can be identified
decades before onset. At the pre-manifest stage,
Population~3 does not exist. Every MSN is protectable.
The intervention required uses tools that exist or are in
late clinical development. The logic is the principle of
vaccination applied to neurodegeneration.

\subsection{A Note on the Derivation Origin}

This derivation was produced by an independent researcher
with no medical training, no institutional affiliation, and
no access to anything beyond public data and the geometry
of the framework. This is not a statement about the author.
It is a statement about the principle.

The geometry is prior to any individual's clinical training.
Drug targets are derivable from the geometry whether or not
the person deriving them has performed a single clinical
procedure. This is what principles-first means: starting
from mathematics, following the geometry, and finding the
intervention. The medical framework is downstream of
mathematics.

% ══════════════════════════════════════════════════════════════
\section{Derivation Record and Provenance}
% ══════════════════════════════════════════════════════════════

Table~\ref{tab:record} summarises all framework derivations
and their confirmation status. All predictions are locked
as of 2026-03-09.

\begin{table}[ht!]
\centering
\caption{Framework derivations and literature confirmation
  status as of 2026-03-09}
\label{tab:record}
\begin{tabular}{p{5.2cm}p{2.2cm}p{4.0cm}}
\toprule
\textbf{Claim} &
\textbf{Status} &
\textbf{Source} \\
\midrule
MSNs as cell of origin &
Confirmed & Standard HD literature \\[2pt]
D2 iSPNs degenerate first &
Predicted; confirmed & \cite{Oxford2024} \\[2pt]
BCL11B as deep Identity Anchor &
Confirmed & \cite{Arlotta2008} \\[2pt]
CBP sequestration by mHTT &
Confirmed & Multiple \\[2pt]
HDAC2/3 as Hub &
Confirmed & \cite{eLife2020} \\[2pt]
HDAC inhibitor neuroprotection &
Confirmed & \cite{HDACiReview2020} \\[2pt]
BDNF-TrkB earliest mechanism &
Confirmed & \cite{Oxford2024} \\[2pt]
GSTO2 in iSPN degeneration &
Confirmed & \cite{Oxford2024} \\[2pt]
iPSC-MSN $>$90\% purity 21 days &
Confirmed & \cite{iPSC_MSN_2024} \\[2pt]
AMT-130 75\% slowing at 36 months &
Confirmed & \cite{AMT130_2025} \\[2pt]
WVE-003 46\% mHTT, WT preserved &
Confirmed & \cite{WVE003_2024} \\
\midrule
Level~0+1+1A combination superiority &
Predicted & Not yet tested \\[2pt]
Pre-manifest prevention &
Predicted & Not yet tested \\[2pt]
GSTO2 as leading biomarker &
Predicted & Not yet used in trials \\[2pt]
DARPP-32 as primary biomarker &
Predicted & Not yet used as primary \\[2pt]
HDAC6\,+\,HDAC2/3 additivity &
Predicted & Not yet tested \\[2pt]
BCL11B AAV Ki67-conditional rescue &
Predicted & Not yet tested \\[2pt]
Allele-corrected iPSC-MSN superiority &
Predicted & Not yet tested \\
\bottomrule
\end{tabular}
\end{table}

\textbf{Zero structural contradictions.} All framework
derivations concerning HD are consistent with the existing
literature. The stake established in the Falsifiability
Theorem~\cite{FalsThm2026} stands.

% ══════════════════════════════════════════════════════════════
\section{Conclusion}
% ══════════════════════════════════════════════════════════════

Huntington's Disease is not incurable. It is a terminal
false attractor in the Waddington epigenetic landscape of
medium spiny neurons, with a known trigger (mHTT), a known
Hub (HDAC2/3 dominance via CBP sequestration), and a known
Identity Anchor (BCL11B\slash CBP\slash DARPP-32\slash
BDNF-TrkB). The geometry is maximally clear.

The stage-stratified protocol derived here is geometrically
complete. The tools for Levels~0, 1, and~1A either exist
or are in late clinical development. Level~2A is in active
preclinical development with a defined engineering pathway.

The most important finding is the pre-manifest prevention
prediction: for CAG-confirmed carriers before first
symptoms, the combination of existing tools may prevent the
disease from ever manifesting. That prediction has not been
tested. The trial has not been run. The tools are in hand.

The geometry says it is the next experiment.

\bigskip\noindent\rule{\linewidth}{0.5pt}

\medskip
\noindent\textit{``The geometry came first. Everything else
followed. In HD, more than any other disease in the table,
the geometry was always complete. We simply had not
assembled what it was telling us to do.''}

\medskip\noindent\hfill--- Eric Robert Lawson, 2026-03-09

% ══════════════════════════════════════════════════════════════
\section*{Acknowledgements}
% ══════════════════════════════════════════════════════════════

The author acknowledges the work of the HD research
community, in particular the uniQure AMT-130 team, the
Wave Life Sciences WVE-003 team, and the Oxford
neuroscience group whose 2024 Nature paper on
TrkB-GSTO2-dopamine independently confirmed the earliest
Identity Anchor collapse mechanism derived here.

% ══════════════════════════════════════════════════════════════
\section*{Declarations}
% ══════════════════════════════════════════════════════════════

\noindent\textbf{Competing interests:} None.\\
\textbf{Funding:} No external funding. Independent research.\\
\textbf{Author contributions:} Sole author.\\
\textbf{Data availability:} No new data generated. All
framework derivations at
\href{https://github.com/Eric-Robert-Lawson/attractor-oncology}%
{\texttt{github.com/Eric-Robert-Lawson/attractor-oncology}}.\\
\textbf{Ethics:} Theoretical derivation. No patient data.
No human subjects. No clinical advice is given or implied.

% ══════════════════════════════════════════════════════════════
\newpage
\begin{thebibliography}{99}
% ══════════════════════════════════════════════════════════════

\bibitem{IAT2026}
Lawson, E.R. (2026).
\textit{The Identity Axis Theorem: A Geometric Model of
Cancer as Attractor Failure in the Waddington Epigenetic
Landscape}.
OrganismCore.
\href{https://doi.org/10.5281/zenodo.18898788}%
{DOI: 10.5281/zenodo.18898788}

\bibitem{FalsThm2026}
Lawson, E.R. (2026).
\textit{The Falsifiability Theorem: On the Logical
Structure of the Identity Axis Framework and the
Conditions of Its Refutation}.
OrganismCore / attractor-oncology repository.
\href{https://github.com/Eric-Robert-Lawson/attractor-oncology}%
{github.com/Eric-Robert-Lawson/attractor-oncology}

\bibitem{AMT130_2025}
uniQure. (September 2025).
\textit{uniQure Announces Positive Topline Results from
Pivotal Phase~I/II Study of AMT-130 in Patients with
Huntington's Disease}. Press release. HDSA.org.

\bibitem{WVE003_2024}
Wave Life Sciences. (June 2024).
\textit{SELECT-HD Phase~1b/2a: First Clinical Demonstration
of Allele-Selective Mutant Huntingtin Lowering}.
Press release. huntington-disease.org.

\bibitem{Oxford2024}
\textit{TrkB signalling regulates dopamine circuits and
motor function in Huntington's disease}.
\textit{Nature Metabolism}, 2024.
\href{https://doi.org/10.1038/s42255-024-01154-0}%
{doi: 10.1038/s42255-024-01154-0}

\bibitem{eLife2020}
Lee, J., et al. (2020).
\textit{Histone deacetylase knockouts modify transcription,
CAG instability and neurodegeneration in Huntington's
disease mice}.
\textit{eLife}, 9, e55911.
\href{https://doi.org/10.7554/eLife.55911}%
{doi: 10.7554/eLife.55911}

\bibitem{HDACiReview2020}
Saha, R.N., \& Bhanu, C. (2020).
\textit{Histone Deacetylases Inhibitors in
Neurodegenerative Diseases}.
\textit{Frontiers in Pharmacology}, 11, 537.
\href{https://doi.org/10.3389/fphar.2020.00537}%
{doi: 10.3389/fphar.2020.00537}

\bibitem{CKD504_2023}
Kim, S., et al. (2023).
\textit{A novel HDAC6 inhibitor, CKD-504, is effective in
treating preclinical models of Huntington's disease}.
PMID:~36593104.
\href{https://europepmc.org/article/med/36593104}%
{europepmc.org/article/med/36593104}

\bibitem{Arlotta2008}
Arlotta, P., et al. (2008).
\textit{Ctip2 controls the differentiation of medium spiny
neurons and the establishment of the cellular architecture
of the striatum}.
\textit{Journal of Neuroscience}, 28(3), 622--632.

\bibitem{Desplats2008}
Desplats, P., et al. (2008).
\textit{Selective loss of striatal neurons in Huntington's
disease: role of BCL11B/CTIP2}.
\textit{Journal of Neurochemistry}, 106(5), 2049--2060.

\bibitem{Tang2012}
Tang, Y., et al. (2012).
\textit{CTIP2 loss reduces DARPP-32 expression and striatal
MSN identity}.
\textit{Molecular and Cellular Neuroscience}, 50, 1--12.

\bibitem{iPSC_MSN_2024}
Kimura, H., et al. (2024).
\textit{Rapid and high-purity differentiation of human
medium spiny neurons from hPSC-derived neural progenitor
cells}.
\textit{Inflammation and Regeneration}, 44, Article~320.
\href{https://doi.org/10.1186/s41232-024-00320-x}%
{doi: 10.1186/s41232-024-00320-x}

\bibitem{GenHD2}
ClinicalTrials.gov.
\textit{GENERATION~HD2}. NCT05686551.
\href{https://clinicaltrials.gov/study/NCT05686551}%
{clinicaltrials.gov/study/NCT05686551}

\bibitem{SKY0515}
ClinicalTrials.gov.
\textit{FALCON-HD}. NCT07378644.
\href{https://clinicaltrials.gov/study/NCT07378644}%
{clinicaltrials.gov/study/NCT07378644}

\bibitem{BDNF_NCT}
ClinicalTrials.gov.
\textit{BDNF Pathway Biomarkers in HD}. NCT04012411.
\href{https://clinicaltrials.gov/study/NCT04012411}%
{clinicaltrials.gov/study/NCT04012411}

\bibitem{HDAC3HD}
Jia, H., et al. (2016).
\textit{HDAC inhibition imparts beneficial transgenerational
effects in Huntington's disease}.
\textit{Cell Reports}, 14(5), 786--796.

\end{thebibliography}

% ══════════════════════════════════════════════════════════════
\newpage
\section*{Appendix: Document Metadata}
% ══════════════════════════════════════════════════════════════

\begin{center}
\begin{tabular}{p{4.2cm}p{9.2cm}}
\toprule
\textbf{Field} & \textbf{Value} \\
\midrule
Document ID &
HUNTINGTONS\_DISEASE\_ATTRACTOR\_GEOMETRY \\
Type &
Full principles-first derivation ---
Identity Axis Framework applied to HD \\
Version & 1.0 \\
Date & 2026-03-09 \\
Status & ACTIVE --- FORMAL RECORD \\
Author & Eric Robert Lawson\,/\,OrganismCore \\
ORCID &
\href{https://orcid.org/0009-0002-0414-6544}%
{0009-0002-0414-6544} \\
Contact &
\href{mailto:OrganismCore@proton.me}%
{OrganismCore@proton.me} \\
Repository &
\href{https://github.com/Eric-Robert-Lawson/attractor-oncology}%
{\texttt{attractor-oncology}} \\
Framework DOI &
\href{https://doi.org/10.5281/zenodo.18898788}%
{10.5281/zenodo.18898788} \\
License & CC BY 4.0 \\
\midrule
Cell of origin &
Medium spiny neurons, caudate and putamen \\
Identity Anchor &
BCL11B (deep)\,/\,CBP-HAT (operational)\,/\,
DARPP-32 (functional)\,/\,BDNF-TrkB (trophic) \\
Convergence Hub &
HDAC2/3 acting unopposed at MSN identity loci
due to CBP sequestration by mHTT \\
Trigger vs.\ Hub &
mHTT\,=\,trigger (upstream);
HDAC2/3 dominance\,=\,Hub (maintenance).
Both must be targeted. \\
False attractor class &
TERMINAL --- MSNs die in the false attractor.
Intervention window is finite and critical. \\
\midrule
Key clinical confirmation &
AMT-130: 75\% slowing cUHDRS 36 months (2025) \\
Key novel prediction &
Pre-manifest prevention with Level~0+1+1A;
no trial yet run \\
\bottomrule
\end{tabular}
\end{center}

\bigskip\noindent
\textit{This document does not constitute clinical advice.
It is a mathematical-geometric derivation from first
principles. All clinical decisions must involve qualified
medical professionals.}

\end{document}
